\documentclass[a4paper,11pt]{article}
\usepackage[utf8]{inputenc}
\usepackage[T1]{fontenc}
\usepackage[french]{babel}
\usepackage{lmodern}
\usepackage{amsmath,amssymb,amsthm}
\usepackage{mathrsfs}
\usepackage{enumitem}
\usepackage{multicol}

\setlist[enumerate]{label=\textbf{\textcolor{Couleur8}{\arabic*.}}, left=1em}

% Si vous voulez aussi modifier les listes enumerate imbriquées :
\setlist[enumerate,2]{label=\textcolor{Couleur8}{\alph*)}}
\setlist[enumerate,3]{label=\textcolor{Couleur8}{\roman*)}}
\setlist[enumerate,4]{label=\textcolor{Couleur8}{\Alph*)}}
\usepackage{graphicx}
\usepackage{xcolor}
\usepackage{tcolorbox}
\usepackage{geometry}
\usepackage{fancyhdr}
\usepackage{titlesec}
\usepackage{cancel}
\usepackage{ifthen}
\usepackage{tikz}
\usetikzlibrary{arrows.meta}
\usepackage{tkz-tab}
\usepackage{eurosym}

% OPTION PROF/ÉLÈVE
% Décommentez UNE des deux lignes suivantes :
\newboolean{prof}\setboolean{prof}{true}   % Version professeur (avec corrections des devoirs)
\newboolean{prof}\setboolean{prof}{true}  % Version élève (sans corrections des devoirs)

\definecolor{Couleur1}{HTML}{4A5D7A} %Th
\definecolor{Couleur2}{HTML}{A5B5D6} %Prop
\definecolor{Couleur3}{HTML}{A8C68A} %Méthode
\definecolor{Couleur6}{HTML}{8A9CC1} %Def
\definecolor{Couleur7}{HTML}{E6B459} %Rq Voc
\definecolor{Couleur8}{HTML}{D36740} %Titres
\definecolor{correction}{HTML}{D36740} % Couleur pour les corrections
\definecolor{coopmathscorr}{HTML}{F15929} % Couleur corrections coopmaths

% Configuration des marges
\geometry{
	top=2cm,
	bottom=2cm,
	inner=1.5cm,
	outer=1.5cm,
}

% Police
\renewcommand{\familydefault}{\sfdefault}

% Compteurs
\newcounter{seancenum}
\newcounter{seancenumD}
\setcounter{seancenum}{1}
\setcounter{seancenumD}{1}

% Configuration des en-têtes et pieds de page
\pagestyle{fancy}
\fancyhf{}
\ifthenelse{\boolean{prof}}{
    \fancyhead[L]{\textcolor{Couleur8}{\textbf{Corrections Automatismes - Classe de seconde - Version Professeur}}}
}{
    \fancyhead[L]{\textcolor{Couleur8}{\textbf{Corrections Devoirs - Séance 11 - Classe de seconde}}}
}
\fancyhead[R]{\textcolor{Couleur8}{\textbf{2025-2026}}}
\fancyfoot[L]{\textcolor{Couleur8}{Mme Bertail et Mme Pinto}}
\fancyfoot[R]{\textcolor{Couleur8}{\thepage}}
\renewcommand{\headrulewidth}{1pt}
\renewcommand{\headrule}{\hbox to\headwidth{\color{Couleur8}\leaders\hrule height \headrulewidth\hfill}}

% Environnements personnalisés
\newtcolorbox{seance}[1]{
    colback=Couleur6!10,
    colframe=Couleur1!65,
    fonttitle=\bfseries\large,
    title=Correction Séance \theseancenum #1,
    left=5mm,
    right=5mm,
    top=3mm,
    bottom=3mm,
    arc=0mm,
    after=\stepcounter{seancenum}
}

\newtcolorbox{seanceD}[1]{
    colback=Couleur6!5,
    colframe=Couleur6!45,
    fonttitle=\bfseries\large,
    title=Correction Devoirs - Séance \theseancenumD #1,
    left=5mm,
    right=5mm,
    top=3mm,
    bottom=3mm,
    arc=0mm,
    after=\stepcounter{seancenumD}
}

\begin{document}

\setcounter{seancenumD}{11}
\newpage
\begin{seanceD}{} %11
\begin{enumerate}
\item  Puisque $f$ est une fonction affine, on a : $f(x)=ax+b$.\\$\bullet$ $b$ est l'ordonnée à l'origine de la droite. On lit $b=-4$.\\$\bullet$ $a$ est le coefficient directeur de la droite.\\
          Il est donné par le déplacement vertical correspondant à un déplacement horizontal d'une unité. On lit $a=3$.\\ On peut en déduire que l'expression de la fonction $f$ est $f(x)={\color[HTML]{f15929}\boldsymbol{3x-4}}$.

\item  $f(x)$ est de la forme $ax+b$, $f$ est donc une fonction affine avec pour coefficient directeur  $a=3$  (positif).\\
             D'où, $f$ est une fonction  croissante, c'est-à-dire que lorsque $x$ augmente, $f(x)$  augmente. \\
             Par conséquent, les valeurs de $f(x)$ sont d'abord négatives puis positives.

\medskip

             Aussi, $3x-3=0 \iff x=1$.

\medskip
 D'où le tableau de signes suivant :\\\begin{tikzpicture}[baseline, scale=0.75]
    \tkzTabInit[lgt=8,deltacl=0.8,espcl=3.5]{ $x$ / 3, $f(x)=3x-3$ / 2}{ $-\infty$, $1$, $+\infty$}
	\tkzTabLine{  , -, z, +}
	\end{tikzpicture}

\item  Pour déterminer le taux d'évolution réciproque, on commence par calculer le coefficient multiplicateur associé :\\Diminuer de $26\,\%$ revient à multiplier par $ 1 - \dfrac{26}{100} = 0{,}74$ 

\medskip
Le coefficient multiplicateur réciproque est donc : \\$CM_R=\dfrac{1}{0{,}74} \approx 1{,}351\,4$.{\color[HTML]{f15929}\\Remarque : Il faut arrondir les valeurs à $10^{-4}$ pour avoir un arrondi en pourcentage à $10^{-2}$.}

\medskip
 Le taux d'évolution réciproque est donc : \\$T_R=CM_R-1=1{,}351\,4-1=0{,}351\,4=35{,}14\,\%$ ce qui correspond à une hausse de $35{,}14\,\%$.

\medskip
Il faut donc appliquer une hausse d'environ $35{,}14\,\%$ pour revenir au niveau de départ.

\item  $t\xrightarrow{\times  4} 4t\xrightarrow{+3}4t+3 \xrightarrow{+3t}4t+3+3t=7t+3$\\Le résultat du programme est donc $7t+3$.

\medskip
Le programme de calcul est équivalent à :	\begin{itemize}
		\item Multiplier par $7$
		\item Ajouter $3$
	\end{itemize}


\item $\dfrac{2x-8}{4}=\dfrac{-8}{5}\\
5( 2x-8 )=4 \times (-8)\\
10x-40=-32\\
10x=-32+40\\
10x=8\\
x=\dfrac{8}{10}=\dfrac{4}{5}$
\\L'équation a donc pour unique solution : ${\color[HTML]{f15929}\boldsymbol{\dfrac{4}{5}}}$.

\end{enumerate}
\end{seanceD}

\end{document}